\scanpage{25}
\mspage{- 217 -}\hfill n\textsuperscript{o} 391

(IV) La racine $\frac{1}{2}(\varepsilon_1+\varepsilon_2+\varepsilon_3+\varepsilon_4+\varepsilon_5-\varepsilon_6-\varepsilon_7+\varepsilon_8)$ est la plus grande des racines de $R''$ pour l'ordre induit par celui considéré sur $R^8$, (n\textsuperscript{o}~10) ; c'est donc la plus grande racine $\tilde\alpha$, (\S~6, n\textsuperscript{o}~6, cor. 2 de la prop. 13) ; on a $\tilde\alpha=\alpha_1+2\alpha_2+2\alpha_3+3\alpha_4+2\alpha_5+\alpha_6$, et on vérifie que son ordre est $h-1=11$. Cette racine est orthogonale à toutes les racines simples $\alpha_i$, $(1\leq i\leq6)$, sauf $\alpha_2$, et l'on a $(\tilde\alpha\mid\alpha_2)=1$.

(V) Le même raisonnement qu'au n\textsuperscript{o}~11 montre que les \uline{coracines} de $R''$ s'obtiennent en remplaçant $\varepsilon_i$ par $\varepsilon_i^*$ dans les expressions des racines de $R''$, et $R''{}^*$ est isomorphe à $R''$. Toutes les racines ayant la même longueur, la \uline{forme bilinéaire canonique} pour $R$ est $(x\mid y)/2h=(x\mid y)/24$ et la longueur des racines est $1/2\sqrt{3}$, (\S~8, n\textsuperscript{o}~5, \uline{Remarque}). Le nombre des racines non orthogonales à une racine donnée est $4h-6=42$ et le nombre $\gamma$ est égal à $h^2=144$.

(VI) Le calcul des poids fondamentaux, formant la base duale de $(\alpha_i^*)$, $(1\leq i\leq6)$, dans $V''$, donne :

\[
\begin{aligned}
\omega_1&=\frac{2}{3}(\varepsilon_8-\varepsilon_7-\varepsilon_6)=\frac{1}{3}(4\alpha_1+3\alpha_2+5\alpha_3+6\alpha_4+4\alpha_5+2\alpha_6)\\
\omega_2&=\frac{1}{2}(\varepsilon_1+\varepsilon_2+\varepsilon_3+\varepsilon_4+\varepsilon_5-\varepsilon_6-\varepsilon_7+\varepsilon_8)=\\
&=\alpha_1+2\alpha_2+2\alpha_3+3\alpha_4+2\alpha_5+\alpha_6=\tilde\alpha\\
\omega_3&=\frac{5}{6}(\varepsilon_8-\varepsilon_7-\varepsilon_6)+\frac{1}{2}(-\varepsilon_1+\varepsilon_2+\varepsilon_3+\varepsilon_4+\varepsilon_5)=\\
&=\frac{1}{3}(5\alpha_1+6\alpha_2+10\alpha_3+12\alpha_4+8\alpha_5+4\alpha_6)\\
\omega_4&=\varepsilon_3+\varepsilon_4+\varepsilon_5-\varepsilon_6-\varepsilon_7+\varepsilon_8=2\alpha_1+3\alpha_2+4\alpha_3+6\alpha_4+4\alpha_5+2\alpha_6\\
\omega_5&=\frac{2}{3}(\varepsilon_8-\varepsilon_7-\varepsilon_6)+\varepsilon_4+\varepsilon_5=\\
&=\frac{1}{3}(4\alpha_1+6\alpha_2+8\alpha_3+12\alpha_4+10\alpha_5+5\alpha_6)\\
\omega_6&=\frac{1}{3}(\varepsilon_8-\varepsilon_7-\varepsilon_6)+\varepsilon_5=\\
&=\frac{1}{3}(2\alpha_1+3\alpha_2+4\alpha_3+6\alpha_4+5\alpha_5+4\alpha_6).
\end{aligned}
\]

% archive sheet number: 24
